\documentclass[11pt]{scrartcl}
\usepackage[sexy]{evan}
\usepackage{quiver}
\usepackage{adjustbox}


\begin{document}

\title{On The GICAR Algebra and The Hausdorff Moment Problem}

\maketitle

\begin{abstract}
\noindent We compute the ordered $K_0$ group associated of the GICAR (Gauge Invariant Canonical Anticommutation Relation) Algebra and use it to solve the Hausdorff moment problem.
\end{abstract}

\tableofcontents

\section{GICAR Algebra and Its Ordered Group}

\begin{definition}[GICAR algebra]
The GICAR algebra is the AF algebra described by the following ($C^*$-algebra inductive limit) Bratteli diagram--Pascal's triangle
% https://q.uiver.app/?q=WzAsOSxbMiwwLCJcXGJ1bGxldCJdLFsxLDEsIlxcYnVsbGV0Il0sWzMsMSwiXFxidWxsZXQiXSxbMCwyLCJcXGJ1bGxldCJdLFsyLDIsIlxcYnVsbGV0Il0sWzQsMiwiXFxidWxsZXQiXSxbMCwzLCJcXHZkb3RzIl0sWzIsMywiXFx2ZG90cyJdLFs0LDMsIlxcdmRvdHMiXSxbMCwxLCIiLDAseyJzdHlsZSI6eyJoZWFkIjp7Im5hbWUiOiJub25lIn19fV0sWzAsMiwiIiwyLHsic3R5bGUiOnsiaGVhZCI6eyJuYW1lIjoibm9uZSJ9fX1dLFsxLDNdLFsxLDQsIiIsMCx7InN0eWxlIjp7ImhlYWQiOnsibmFtZSI6Im5vbmUifX19XSxbMiw0LCIiLDIseyJzdHlsZSI6eyJoZWFkIjp7Im5hbWUiOiJub25lIn19fV0sWzIsNSwiIiwyLHsic3R5bGUiOnsiaGVhZCI6eyJuYW1lIjoibm9uZSJ9fX1dXQ==
\[\begin{tikzcd}
	&& \bullet \\
	& \bullet && \bullet \\
	\bullet && \bullet && \bullet \\
	\vdots && \vdots && \vdots
	\arrow[no head, from=1-3, to=2-2]
	\arrow[no head, from=1-3, to=2-4]
	\arrow[from=2-2, to=3-1]
	\arrow[no head, from=2-2, to=3-3]
	\arrow[no head, from=2-4, to=3-3]
	\arrow[no head, from=2-4, to=3-5]
\end{tikzcd}\]
\end{definition}

\noindent Let's first compute the $K_0$ group of the GICAR algebra.

\begin{theorem}[$K_0$ of GICAR]\label{K_0 of GICAR}
Continuity of $K_0$ (Theorem 6.3.2 in \cite{rørdam_larsen_laustsen_2010}) gives us that the $K_0$ group of the GICAR algebra is described by the following (group inductive limit) Bratteli diagram
% https://q.uiver.app/?q=WzAsOSxbMiwwLCJcXG1hdGhiYntafSJdLFsxLDEsIlxcbWF0aGJie1p9Il0sWzMsMSwiXFxtYXRoYmJ7Wn0iXSxbMCwyLCJcXG1hdGhiYntafSJdLFsyLDIsIlxcbWF0aGJie1p9Il0sWzQsMiwiXFxtYXRoYmJ7Wn0iXSxbMCwzLCJcXHZkb3RzIl0sWzIsMywiXFx2ZG90cyJdLFs0LDMsIlxcdmRvdHMiXSxbMCwxLCIiLDAseyJzdHlsZSI6eyJoZWFkIjp7Im5hbWUiOiJub25lIn19fV0sWzAsMiwiIiwyLHsic3R5bGUiOnsiaGVhZCI6eyJuYW1lIjoibm9uZSJ9fX1dLFsxLDMsIiIsMCx7InN0eWxlIjp7ImhlYWQiOnsibmFtZSI6Im5vbmUifX19XSxbMSw0LCIiLDAseyJzdHlsZSI6eyJoZWFkIjp7Im5hbWUiOiJub25lIn19fV0sWzIsNCwiIiwyLHsic3R5bGUiOnsiaGVhZCI6eyJuYW1lIjoibm9uZSJ9fX1dLFsyLDUsIiIsMix7InN0eWxlIjp7ImhlYWQiOnsibmFtZSI6Im5vbmUifX19XV0=
\[\begin{tikzcd}
	&& {\mathbb{Z}} \\
	& {\mathbb{Z}} && {\mathbb{Z}} \\
	{\mathbb{Z}} && {\mathbb{Z}} && {\mathbb{Z}} \\
	\vdots && \vdots && \vdots
	\arrow[no head, from=1-3, to=2-2]
	\arrow[no head, from=1-3, to=2-4]
	\arrow[no head, from=2-2, to=3-1]
	\arrow[no head, from=2-2, to=3-3]
	\arrow[no head, from=2-4, to=3-3]
	\arrow[no head, from=2-4, to=3-5]
\end{tikzcd}\]
Note that the edges represent the group homomorphism taking $1\mapsto 1.$
This group is $\ZZ[x]$.
\end{theorem}
\begin{proof}
$K_0$ of the GICAR algebra is isomorphic to the group generated by the following Bratteli diagram (*)
% https://q.uiver.app/?q=WzAsOSxbMiwwLCJcXGxhbmdsZSAxIFxccmFuZ2xlIl0sWzEsMSwiXFxsYW5nbGUgeCBcXHJhbmdsZSJdLFszLDEsIlxcbGFuZ2xlIDEgLSB4IFxccmFuZ2xlIl0sWzAsMiwiXFxsYW5nbGUgeF4yIFxccmFuZ2xlIl0sWzIsMiwiXFxsYW5nbGUgeCgxIC0geCkgXFxyYW5nbGUiXSxbNCwyLCJcXGxhbmdsZSAoMSAtIHgpXjIgXFxyYW5nbGUiXSxbMCwzLCJcXHZkb3RzIl0sWzIsMywiXFx2ZG90cyJdLFs0LDMsIlxcdmRvdHMiXSxbMCwxLCIiLDAseyJzdHlsZSI6eyJoZWFkIjp7Im5hbWUiOiJub25lIn19fV0sWzAsMiwiIiwyLHsic3R5bGUiOnsiaGVhZCI6eyJuYW1lIjoibm9uZSJ9fX1dLFsxLDMsIiIsMCx7InN0eWxlIjp7ImhlYWQiOnsibmFtZSI6Im5vbmUifX19XSxbMSw0LCIiLDAseyJzdHlsZSI6eyJoZWFkIjp7Im5hbWUiOiJub25lIn19fV0sWzIsNCwiIiwyLHsic3R5bGUiOnsiaGVhZCI6eyJuYW1lIjoibm9uZSJ9fX1dLFsyLDUsIiIsMix7InN0eWxlIjp7ImhlYWQiOnsibmFtZSI6Im5vbmUifX19XV0=
\[\begin{tikzcd}
	&& {\langle 1 \rangle} \\
	& {\langle x \rangle} && {\langle 1 - x \rangle} \\
	{\langle x^2 \rangle} && {\langle x(1 - x) \rangle} && {\langle (1 - x)^2 \rangle} \\
	\vdots && \vdots && \vdots
	\arrow[no head, from=1-3, to=2-2]
	\arrow[no head, from=1-3, to=2-4]
	\arrow[no head, from=2-2, to=3-1]
	\arrow[no head, from=2-2, to=3-3]
	\arrow[no head, from=2-4, to=3-3]
	\arrow[no head, from=2-4, to=3-5]
\end{tikzcd}\]
where each entry is a subgroup of $\ZZ[x].$
For proof, consider the following diagram, where the red arrows map generator to generator, where the generator for $\ZZ$ is $1$
\\

% https://q.uiver.app/?q=WzAsMTgsWzIsMCwiXFxtYXRoYmJ7Wn0iXSxbMSwxLCJcXG1hdGhiYntafSJdLFszLDEsIlxcbWF0aGJie1p9Il0sWzAsMiwiXFxtYXRoYmJ7Wn0iXSxbMiwyLCJcXG1hdGhiYntafSJdLFs0LDIsIlxcbWF0aGJie1p9Il0sWzAsMywiXFx2ZG90cyJdLFsyLDMsIlxcdmRvdHMiXSxbNCwzLCJcXHZkb3RzIl0sWzgsMCwiXFxsYW5nbGUgMSBcXHJhbmdsZSJdLFs3LDEsIlxcbGFuZ2xlIHggXFxyYW5nbGUiXSxbOSwxLCJcXGxhbmdsZSAxIC0geCBcXHJhbmdsZSJdLFs2LDIsIlxcbGFuZ2xlIHheMiBcXHJhbmdsZSJdLFs4LDIsIlxcbGFuZ2xlIHgoMSAtIHgpIFxccmFuZ2xlIl0sWzEwLDIsIlxcbGFuZ2xlICgxIC0geCleMiBcXHJhbmdsZSJdLFs2LDMsIlxcdmRvdHMiXSxbOCwzLCJcXHZkb3RzIl0sWzEwLDMsIlxcdmRvdHMiXSxbMCwxLCIiLDAseyJzdHlsZSI6eyJoZWFkIjp7Im5hbWUiOiJub25lIn19fV0sWzAsMiwiIiwyLHsic3R5bGUiOnsiaGVhZCI6eyJuYW1lIjoibm9uZSJ9fX1dLFsxLDMsIiIsMCx7InN0eWxlIjp7ImhlYWQiOnsibmFtZSI6Im5vbmUifX19XSxbMSw0LCIiLDAseyJzdHlsZSI6eyJoZWFkIjp7Im5hbWUiOiJub25lIn19fV0sWzIsNCwiIiwyLHsic3R5bGUiOnsiaGVhZCI6eyJuYW1lIjoibm9uZSJ9fX1dLFsyLDUsIiIsMix7InN0eWxlIjp7ImhlYWQiOnsibmFtZSI6Im5vbmUifX19XSxbOSwxMCwiIiwwLHsic3R5bGUiOnsiaGVhZCI6eyJuYW1lIjoibm9uZSJ9fX1dLFs5LDExLCIiLDIseyJzdHlsZSI6eyJoZWFkIjp7Im5hbWUiOiJub25lIn19fV0sWzEwLDEyLCIiLDAseyJzdHlsZSI6eyJoZWFkIjp7Im5hbWUiOiJub25lIn19fV0sWzEwLDEzLCIiLDAseyJzdHlsZSI6eyJoZWFkIjp7Im5hbWUiOiJub25lIn19fV0sWzExLDEzLCIiLDIseyJzdHlsZSI6eyJoZWFkIjp7Im5hbWUiOiJub25lIn19fV0sWzExLDE0LCIiLDIseyJzdHlsZSI6eyJoZWFkIjp7Im5hbWUiOiJub25lIn19fV0sWzAsOSwiIiwxLHsiY3VydmUiOi0zLCJjb2xvdXIiOlswLDYwLDYwXSwic3R5bGUiOnsidGFpbCI6eyJuYW1lIjoiYXJyb3doZWFkIn19fV0sWzIsMTEsIiIsMSx7ImN1cnZlIjotMywiY29sb3VyIjpbMCw2MCw2MF0sInN0eWxlIjp7InRhaWwiOnsibmFtZSI6ImFycm93aGVhZCJ9fX1dLFsxLDEwLCIiLDEseyJjdXJ2ZSI6LTMsImNvbG91ciI6WzAsNjAsNjBdLCJzdHlsZSI6eyJ0YWlsIjp7Im5hbWUiOiJhcnJvd2hlYWQifX19XSxbMywxMiwiIiwxLHsiY3VydmUiOi0zLCJjb2xvdXIiOlswLDYwLDYwXSwic3R5bGUiOnsidGFpbCI6eyJuYW1lIjoiYXJyb3doZWFkIn19fV0sWzQsMTMsIiIsMSx7ImN1cnZlIjotMywiY29sb3VyIjpbMCw2MCw2MF0sInN0eWxlIjp7InRhaWwiOnsibmFtZSI6ImFycm93aGVhZCJ9fX1dLFs1LDE0LCIiLDEseyJjdXJ2ZSI6LTMsImNvbG91ciI6WzAsNjAsNjBdLCJzdHlsZSI6eyJ0YWlsIjp7Im5hbWUiOiJhcnJvd2hlYWQifX19XV0=
\adjustbox{scale=0.75,center}{%
\begin{tikzcd}
	&& {\mathbb{Z}} &&&&&& {\langle 1 \rangle} \\
	& {\mathbb{Z}} && {\mathbb{Z}} &&&& {\langle x \rangle} && {\langle 1 - x \rangle} \\
	{\mathbb{Z}} && {\mathbb{Z}} && {\mathbb{Z}} && {\langle x^2 \rangle} && {\langle x(1 - x) \rangle} && {\langle (1 - x)^2 \rangle} \\
	\vdots && \vdots && \vdots && \vdots && \vdots && \vdots
	\arrow[no head, from=1-3, to=2-2]
	\arrow[no head, from=1-3, to=2-4]
	\arrow[no head, from=2-2, to=3-1]
	\arrow[no head, from=2-2, to=3-3]
	\arrow[no head, from=2-4, to=3-3]
	\arrow[no head, from=2-4, to=3-5]
	\arrow[no head, from=1-9, to=2-8]
	\arrow[no head, from=1-9, to=2-10]
	\arrow[no head, from=2-8, to=3-7]
	\arrow[no head, from=2-8, to=3-9]
	\arrow[no head, from=2-10, to=3-9]
	\arrow[no head, from=2-10, to=3-11]
	\arrow[color={rgb,255:red,214;green,92;blue,92}, curve={height=-18pt}, tail reversed, from=1-3, to=1-9]
	\arrow[color={rgb,255:red,214;green,92;blue,92}, curve={height=-18pt}, tail reversed, from=2-4, to=2-10]
	\arrow[color={rgb,255:red,214;green,92;blue,92}, curve={height=-18pt}, tail reversed, from=2-2, to=2-8]
	\arrow[color={rgb,255:red,214;green,92;blue,92}, curve={height=-18pt}, tail reversed, from=3-1, to=3-7]
	\arrow[color={rgb,255:red,214;green,92;blue,92}, curve={height=-18pt}, tail reversed, from=3-3, to=3-9]
	\arrow[color={rgb,255:red,214;green,92;blue,92}, curve={height=-18pt}, tail reversed, from=3-5, to=3-11]
\end{tikzcd}
}
Observe that in each triangle of (*), we have $A = B + C$ (**).
% https://q.uiver.app/?q=WzAsNixbMSwxLCJcXGxhbmdsZSBBIFxccmFuZ2xlIl0sWzAsMiwiXFxsYW5nbGUgQlxccmFuZ2xlIl0sWzIsMiwiXFxsYW5nbGUgQyBcXHJhbmdsZSJdLFswLDMsIlxcdmRvdHMiXSxbMiwzLCJcXHZkb3RzIl0sWzEsMCwiXFx2ZG90cyJdLFswLDEsIiIsMSx7InN0eWxlIjp7ImhlYWQiOnsibmFtZSI6Im5vbmUifX19XSxbMCwyLCIiLDEseyJzdHlsZSI6eyJoZWFkIjp7Im5hbWUiOiJub25lIn19fV1d
\[\begin{tikzcd}
	& \vdots \\
	& {\langle A \rangle} \\
	{\langle B\rangle} && {\langle C \rangle} \\
	\vdots && \vdots
	\arrow[no head, from=2-2, to=3-1]
	\arrow[no head, from=2-2, to=3-3]
\end{tikzcd}\]
Therefore, we can treat the group direct sums as simply addition: ie. we can treat $\langle x^2\rangle\oplus \langle x(1 - x)\rangle\oplus \langle (1 - x)^2\rangle$ as $\langle a_1x^2 + a_2x(1 - x) + a_3(1 - x)^2\ |\ a_i\in \ZZ\rangle.$
Moreover, (**) means that $\langle a_1x^n + a_2x^{n - 1}(1 - x) + ... + a_n(1 - x)^n\ |\ a_i\in \ZZ\rangle = \ZZ[x]_{\le n}$ (degree $\le$ $n$) because $1, x, ..., x^n$ belong to this subgroup and the max degree polynomial in the subgroup is degree $n.$
Thus, the $K_0$ group we are computing is the inductive limit 
% https://q.uiver.app/?q=WzAsNCxbMCwwLCJcXG1hdGhiYntafVt4XV97XFxsZSAwfSJdLFsyLDAsIlxcbWF0aGJie1p9W3hdX3tcXGxlIDF9Il0sWzQsMCwiXFxtYXRoYmJ7Wn1beF1fe1xcbGUgMn0iXSxbNiwwLCJcXGNkb3RzIl0sWzAsMSwiXFxpb3RhXzAiXSxbMSwyLCJcXGlvdGFfMSJdLFsyLDMsIlxcaW90YV8yIl1d
\[\begin{tikzcd}
	{\mathbb{Z}[x]_{\le 0}} && {\mathbb{Z}[x]_{\le 1}} && {\mathbb{Z}[x]_{\le 2}} && \cdots
	\arrow["{\iota_0}", from=1-1, to=1-3]
	\arrow["{\iota_1}", from=1-3, to=1-5]
	\arrow["{\iota_2}", from=1-5, to=1-7]
\end{tikzcd}\]
where the connecting maps are inclusions by (**).
Hence, this group must be $\ZZ[x].$
\end{proof}

\noindent Now let's take a look at the ordred $K_0$ group.
We first present a result of Halperin in \cite{Czarnecki1964}.

\begin{lemma}\label{halperin}
Let $p(x)\in \RR[x]$ be non-negative on the interval $[a, b].$
Then, $p(x)$ is the sum of the following forms: $(x - a)q(x)^2,$ $(b - x)q(x)^2,$ $q(x)^2$ where $q(x)\in \RR[x].$
\end{lemma}
\begin{proof}
We proceed by induction on the degree of $p(x).$
If $p(x) = c,$ then we can write $c = (\sqrt{c})^2.$
Now suppose $p(x)$ is a degree $n\ge 1$ polynomial.
Taking $m = \min_{x\in [a, b]} p(x),$ we have $p(x) = p(x) - m + (\sqrt{m})^2,$ so it suffices to prove the hypothesis for $p(x) - m.$
Since $m = \min_{x\in [a, b]} p(x),$ we can factor $p(x) - m = (x - r)q(x)$ for some $r\in [a, b].$
Here we consider two cases.
\begin{enumerate}
    \item If $r\in (a, b),$ then $q(r) = 0$ since $x - r$ changing sign as $x$ passes $r$ forces $q(x)$ to change sign with it.
    Hence, $p(x) - m = (x - r)q(x) = (x - r)^2s(x).$
    Since $\deg s(x) < n,$ we can invoke inductive hypothesis and are done.
    \item Suppose $r = a$ or $r = b.$
    Without loss of generality, suppose $r = a,$ for the $r = b$ case is symmetric.
    We have $p(x) - m = (x - a)q(x).$
    Since $\deg q(x) < n,$ we know $q(x)$ is the sum of the forms $(x - a)s(x)^2,$ $(b - x)s(x)^2,$ $s(x)^2.$
    It suffices to prove that our inductive hypothesis holds for $p(x) - m$ when $q(x) = (b - x)s(x)^2.$
    Observe that 
    \begin{align*}
        (x - a)(b - x) &= (x - a)(b - x)\frac{(x - a + b - x)}{b - a}\\
        &= \frac{1}{b - a}[(b - x)(x - a)^2 + (x - a)(b - x)^2].
    \end{align*}
    Since $\frac{1}{b - a} > 0,$ we can put $\frac{1}{\sqrt{b - a}}$ into each of the squares.
    Then, $p(x) - m$ satisfies our inductive hypothesis.
\end{enumerate}
\end{proof}

\begin{lemma}\label{basis for positive cone}
A polynomial $p(x)\in \ZZ[x]$ is non-negative on $\{0, 1\}$ and strictly positive on $(0, 1)$ if and only if it can be written as 
$$
    p(x) = \sum_k a_kx^{e_k}(1 - x)^{f_k}
$$
where $a_k\in \ZZ$ and $a_k > 0.$
\end{lemma}
\begin{proof}
The $\impliedby$ direction is easy.
Let's consider the $\implies$ direction.
Suppose $p(x)\in \ZZ[x]$ is non-negative on $\{0, 1\}$ and strictly positive on $(0, 1).$
We will proceed by induction on the degree of $p(x)$.
It's clear $p(x)\neq 0.$
If $p(x)$ is degree $0,$ $p(x) = c > 0.$
Note that $c = cx + c(1 - x).$
Suppose that the hypothesis holds when $\deg p(x)\le n.$
By Lemma \ref{halperin}, we are done.
\end{proof}

\begin{remark}
One can see that we can swap $\ZZ$ with $\RR$ without any problems.
\end{remark}

\begin{theorem}[Ordered $K_0$ of GICAR]
The ordered $K_0$ group of GICAR is $\ZZ[x],$ where $\ZZ[x]^+$ consists of $0$ and elements of $\ZZ[x]$ that are non-negative on $\{0, 1\}$ and strictly positive on $(0, 1).$
\end{theorem}
\begin{proof}
A direct consequence of continuity of $K_0$ (Theorem 6.3.2 in \cite{rørdam_larsen_laustsen_2010}) and Lemma \ref{basis for positive cone}.    
\end{proof}

\section{Hausdorff Moment Problem}

The Hausdorff moment problem asks for necessary and sufficient conditions for a sequence $a_0, a_1, ...$ to satisfy 
$$
    a_k = \int_0^1 x^k\ d\mu(x)
$$
for some positive measure $\mu.$
We will henceforth suppress the integration bounds.

\begin{theorem}[Necessary condition]\label{necessary condition}
Let $a_0, a_1, ...$ satisfy 
$$
    a_k = \int x^k\ d\mu(x)
$$
for some positive measure $\mu.$
Then, the Pascal triangle generated by writing $a_0, a_1, ...$ along the left diagonal must have non-negative entries, ie. every entry in the following diagram must be non-negative.

% https://q.uiver.app/?q=WzAsMTAsWzMsMCwiYV8wIl0sWzIsMSwiYV8xIl0sWzEsMiwiYV8yIl0sWzQsMSwiYV8wIC0gYV8xIl0sWzMsMiwiYV8xIC0gYV8yIl0sWzAsMywiYV8zIl0sWzIsMywiYV8yIC0gYV8zIl0sWzQsMywiXFxkZG90cyJdLFs1LDIsIlxcZGRvdHMiXSxbNiwzLCJcXGRkb3RzIl0sWzAsMSwiIiwwLHsic3R5bGUiOnsiaGVhZCI6eyJuYW1lIjoibm9uZSJ9fX1dLFswLDMsIiIsMix7InN0eWxlIjp7ImhlYWQiOnsibmFtZSI6Im5vbmUifX19XSxbMSwyLCIiLDAseyJzdHlsZSI6eyJoZWFkIjp7Im5hbWUiOiJub25lIn19fV0sWzEsNCwiIiwwLHsic3R5bGUiOnsiaGVhZCI6eyJuYW1lIjoibm9uZSJ9fX1dLFszLDQsIiIsMix7InN0eWxlIjp7ImhlYWQiOnsibmFtZSI6Im5vbmUifX19XSxbMyw4LCIiLDIseyJzdHlsZSI6eyJoZWFkIjp7Im5hbWUiOiJub25lIn19fV0sWzIsNSwiIiwwLHsic3R5bGUiOnsiaGVhZCI6eyJuYW1lIjoibm9uZSJ9fX1dLFsyLDYsIiIsMCx7InN0eWxlIjp7ImhlYWQiOnsibmFtZSI6Im5vbmUifX19XSxbNCw2LCIiLDAseyJzdHlsZSI6eyJoZWFkIjp7Im5hbWUiOiJub25lIn19fV0sWzQsNywiIiwwLHsic3R5bGUiOnsiaGVhZCI6eyJuYW1lIjoibm9uZSJ9fX1dLFs4LDcsIiIsMix7InN0eWxlIjp7ImhlYWQiOnsibmFtZSI6Im5vbmUifX19XSxbOCw5LCIiLDIseyJzdHlsZSI6eyJoZWFkIjp7Im5hbWUiOiJub25lIn19fV1d&macro_url=%5Cdef%5Crddots%231%7B%5Ccdot%5E%7B%5Ccdot%5E%7B%5Ccdot%5E%7B%231%7D%7D%7D%7D
\[\begin{tikzcd}
	&&& {a_0} \\
	&& {a_1} && {a_0 - a_1} \\
	& {a_2} && {a_1 - a_2} && \ddots \\
	{a_3} && {a_2 - a_3} && \ddots && \ddots
	\arrow[no head, from=1-4, to=2-3]
	\arrow[no head, from=1-4, to=2-5]
	\arrow[no head, from=2-3, to=3-2]
	\arrow[no head, from=2-3, to=3-4]
	\arrow[no head, from=2-5, to=3-4]
	\arrow[no head, from=2-5, to=3-6]
	\arrow[no head, from=3-2, to=4-1]
	\arrow[no head, from=3-2, to=4-3]
	\arrow[no head, from=3-4, to=4-3]
	\arrow[no head, from=3-4, to=4-5]
	\arrow[no head, from=3-6, to=4-5]
	\arrow[no head, from=3-6, to=4-7]
\end{tikzcd}\]
\end{theorem}
\begin{proof}
Integral of positive function with respect to positive measure is positive.
Thus, our condition is indeed a necessary condition because by replacing the values in the Pascal triangle with their integral counterparts, we get the following diagram.
\\

% https://q.uiver.app/?q=WzAsOSxbMywwLCJcXGludCAxXFwgZFxcbXUoeCkiXSxbMiwxLCJcXGludCB4XFwgZFxcbXUoeCkiXSxbNCwxLCJcXGludCAxIC0geFxcIGRcXG11KHgpIl0sWzMsMiwiXFxpbnQgeCgxIC0geClcXCBkXFxtdSh4KSJdLFswLDMsIlxcaW50IHheM1xcIGRcXG11KHgpIl0sWzEsMiwiXFxpbnQgeF4yXFwgZFxcbXUoeCkiXSxbMiwzLCJcXGludCB4XjIoMSAtIHgpXFwgZFxcbXUoeCkiXSxbNSwyXSxbNCwzXSxbMCwxLCIiLDAseyJzdHlsZSI6eyJoZWFkIjp7Im5hbWUiOiJub25lIn19fV0sWzAsMiwiIiwyLHsic3R5bGUiOnsiaGVhZCI6eyJuYW1lIjoibm9uZSJ9fX1dLFsxLDMsIiIsMCx7InN0eWxlIjp7ImhlYWQiOnsibmFtZSI6Im5vbmUifX19XSxbMiwzLCIiLDIseyJzdHlsZSI6eyJoZWFkIjp7Im5hbWUiOiJub25lIn19fV0sWzEsNSwiIiwwLHsic3R5bGUiOnsiaGVhZCI6eyJuYW1lIjoibm9uZSJ9fX1dLFs1LDQsIiIsMCx7InN0eWxlIjp7ImhlYWQiOnsibmFtZSI6Im5vbmUifX19XSxbNSw2LCIiLDAseyJzdHlsZSI6eyJoZWFkIjp7Im5hbWUiOiJub25lIn19fV0sWzMsNiwiIiwwLHsic3R5bGUiOnsiaGVhZCI6eyJuYW1lIjoibm9uZSJ9fX1dLFsyLDcsIlxcZGRvdHMiLDIseyJvZmZzZXQiOi01LCJzdHlsZSI6eyJib2R5Ijp7Im5hbWUiOiJub25lIn0sImhlYWQiOnsibmFtZSI6Im5vbmUifX19XSxbMyw4LCJcXGRkb3RzIiwyLHsib2Zmc2V0IjotNSwic3R5bGUiOnsiYm9keSI6eyJuYW1lIjoibm9uZSJ9LCJoZWFkIjp7Im5hbWUiOiJub25lIn19fV1d&macro_url=%5Cdef%5Crddots%231%7B%5Ccdot%5E%7B%5Ccdot%5E%7B%5Ccdot%5E%7B%231%7D%7D%7D%7D
\adjustbox{scale=0.8,center}{%
\begin{tikzcd}
	&&& {\int 1\ d\mu(x)} \\
	&& {\int x\ d\mu(x)} && {\int 1 - x\ d\mu(x)} \\
	& {\int x^2\ d\mu(x)} && {\int x(1 - x)\ d\mu(x)} && {} \\
	{\int x^3\ d\mu(x)} && {\int x^2(1 - x)\ d\mu(x)} && {}
	\arrow[no head, from=1-4, to=2-3]
	\arrow[no head, from=1-4, to=2-5]
	\arrow[no head, from=2-3, to=3-4]
	\arrow[no head, from=2-5, to=3-4]
	\arrow[no head, from=2-3, to=3-2]
	\arrow[no head, from=3-2, to=4-1]
	\arrow[no head, from=3-2, to=4-3]
	\arrow[no head, from=3-4, to=4-3]
	\arrow["\ddots"', shift left=5, draw=none, from=2-5, to=3-6]
	\arrow["\ddots"', shift left=5, draw=none, from=3-4, to=4-5]
\end{tikzcd}
}
\end{proof}

\begin{lemma}\label{weierstrass positive}
Functions $f\in C([0, 1])$ such that $f\ge 0$ can be uniformly approximated by polynomials that are strictly positive on $[0, 1].$
\end{lemma}
\begin{proof}
Let $f\in C([0, 1])$ where $f\ge 0$ and $\epsilon > 0$ be arbitrary.
Observe that $\sqrt{f}\in C([0, 1])$ so for any $\delta > 0,$ by the Weierstrass theorem, there exist a polynomial $p$ such that $\sup_{x\in [0, 1]} |\sqrt{f(x)} - p(x)| < \delta.$
Let $M = \sup_{x\in [0, 1]} |\sqrt{f(x)}|.$
We have 
\begin{align*}
    |f(x) - p(x)^2| &= |\sqrt{f(x)} - p(x)|\cdot |\sqrt{f(x)} - \sqrt{f(x)} + \sqrt{f(x)} + p(x)|\\
    &\le \delta\cdot (|p(x) - \sqrt{f(x)}| + 2|\sqrt{f(x)}|)\\
    &\le \delta\cdot (\delta + 2M)\\
    &\rightarrow 0.
\end{align*}
This means that there exist a polynomial $q\ge 0$ such that $\sup_{x\in [0, 1]} |f(x) - q(x)| < \frac{\epsilon}{2}.$
Then, take $q(x) + \frac{\epsilon}{2}$ to be our strictly positive approximating polynomial: $\sup_{x\in [0, 1]} |f(x) - (q(x) + \frac{\epsilon}{2})| < \epsilon.$
We are done.
\end{proof}

\begin{theorem}
The necessary condition is also a sufficient condition.
\end{theorem}
\begin{proof}
Suppose the sequence $a_0, a_1, ...$ satisfies the necessary condition in Theorem \ref{necessary condition}.
By the Weierstrass theorem we can define a bounded (***) linear functional $\psi: C([0, 1])\rightarrow \RR$ by setting $\psi(x^k) = a_k.$
We claim that $\psi$ is a positive linear functional.
Let $f\in C([0, 1])$ be positive.
Lemma \ref{weierstrass positive} tells us that there exist a sequence of polynomials $p_n,$ each non-negative on $\{0, 1\}$ and strictly positive on $(0, 1),$ that converge uniformly to $f.$
Because $a_1, a_2, ...$ satisfies the necessary condition, we know $\psi(x^e(1 - x)^f)\ge 0$ for all valid $e, f.$
Then, by Lemma \ref{basis for positive cone} and continuity of $\psi,$ we know that $\psi(f)\ge 0.$
Finally, we can conclude by using the RMK theorem that there exists a Randon (hence positive) measure $\mu$ such that for all $f\in C([0, 1])$
$$
    \psi(f) = \int f(x)\ d\mu(x).
$$
We are done.
\end{proof}

\begin{remark}
We are not sure how to prove (***).
\end{remark}

\bibliographystyle{plain}
\bibliography{refs}

\end{document}